\section{Introduction}\label{sec:introduction}

\subsection{Problem context}
Urban mobility data arrives continuously and has immediate operational relevance: each trip contributes pickup location, distance, fare, and zone-demand observations. An Intelligent Transportation System (ITS) -- a transport platform that uses real-time data to coordinate dispatch, analytics, and planning -- requires more than delayed batch analysis; operators also need near-real-time aggregates and a way to modify analytical rules when data-quality or business requirements evolve. ``Near-real-time'' here means that results must reflect the most recent events within seconds to tens of seconds rather than after a scheduled batch run.

In batch processing, the system collects a body of data and computes results on a schedule, so an analytical view commonly represents the past after a planned delay. Stream processing instead treats each newly emitted record as an \emph{event} -- a message that can update a served result the moment it arrives. A key abstraction that makes this possible at the analytical layer is the \emph{materialized view} (MV): a query result that the system maintains incrementally as input changes, so readers query the maintained aggregate rather than recomputing from raw records. This distinction matters for ITS: a changing zone-demand aggregate is operationally useful only if it can be observed soon enough and if its definition can be modified safely while it is being served.

The lakehouse model is an architecture that combines the flexibility of a \emph{data lake} (open files on inexpensive object storage) with the managed table semantics of a \emph{data warehouse}, and was formally proposed in prior research~\cite{armbrust2021lakehouse}. Stream processing adds further components to the operational path: a \emph{broker} (intermediate layer that stores and routes messages), a computation engine, a \emph{state store} (which holds intermediate state across events), a \emph{sink} (output that writes to a downstream system), and an observability layer. The Ursa paper \parencite{merli2025ursa} identifies the cost and complexity that arise when streaming and lakehouse infrastructure are separated, and proposes a Kafka-compatible lakehouse-native engine. At the analytical application layer, a practical question remains widely unanswered: can the definition of a serving aggregate be changed without observed query failures during the transition?

\subsection{Motivation}
Continux studies the question above in a setting that is reproducible in a laboratory: a lightweight Kubernetes cluster, the public NYC TLC taxi dataset, and standard open-source components. The chosen Kubernetes distribution is K3s -- a lightweight Kubernetes distribution maintained by SUSE that can run on resource-constrained machines -- which allows the project to exercise core Kubernetes concepts without large-scale cloud infrastructure. The ITS motivation also relates to the United Nations Sustainable Development Goals (SDG): SDG 8 through efficient mobility-supporting digital infrastructure, SDG 9 through data infrastructure innovation, and SDG 11 through observable urban mobility analytics \parencite{sdg8,sdg9,sdg11}. These connections motivate the application domain; they do not constitute a quantified SDG impact evaluation.

Ursa primarily addresses streaming-engine and ingestion architecture. Continux does not reproduce Ursa's cost benchmark; it instead examines the \emph{operational layer} above ingestion, comprising GitOps-managed deployment, observability dashboards, and a RisingWave materialized-view name swap used to evolve public analytical logic, which is detailed in Section~\ref{sec:background}.

\subsection{Objectives and scope}
The study has four coordinated objectives that move from design to evaluation:

\begin{itemize}
\item Design a streaming lakehouse pipeline from NYC TLC data to Apache Iceberg objects on object storage.
\item Deploy the pipeline on a three-server HA K3s cluster with separate data, control/observability, and quorum-only roles.
\item Execute a Blue/Green cutover at the public materialized-view query layer while observing duration, query errors, and workload status.
\item Evaluate the artifact strictly from collected evidence and dashboard captures, without extending conclusions beyond what was measured.
\end{itemize}

The processed source is the public NYC TLC Yellow Taxi dataset for \texttt{2026-03}, joined with the Taxi Zone Lookup table. The Parquet input -- a columnar compressed format optimized for analytics -- is converted to \texttt{3,952,451} JSONL records. The reported experiment deliberately stops the Vector replay process after the public MV reaches \texttt{986} trips; the final MV therefore reflects a sample rather than the full converted file. Cutover conclusions apply to the public query name \texttt{mv\_zone\_stats}; the study does not claim that the downstream Iceberg sink automatically migrates to green logic as a result of the materialized-view name swap.

\subsection{Research questions}
\begin{table*}[t!]
\centering
\caption{Research questions and available evidence}
\label{tab:rqs-en}
\begin{tabularx}{\textwidth}{@{}p{0.08\textwidth}Y Y Y@{}}
\toprule
\textbf{ID} & \textbf{Question} & \textbf{Evidence} & \textbf{Bounded conclusion}\\
\midrule
RQ1 & Can public analytical logic be switched without observed query errors during the swap window? & Runbook duration \texttt{0.145226 s}; query errors \texttt{0}. & Yes, at query layer in the observed window.\\
RQ2 & Does the pipeline generate observable lakehouse output? & MinIO lists Parquet data, delete objects, and metadata. & Yes, after replay.\\
RQ3 & Can the constrained cluster run the designed scenario? & \texttt{3/3} nodes, \texttt{5/5} PVCs, \texttt{23/23} workloads. & Yes, for this scenario.\\
RQ4 & Has end-to-end exactly-once behavior been verified? & No offset/duplicate/loss reconciliation. & Not established.\\
\bottomrule
\end{tabularx}
\end{table*}

\subsection{Contributions and organization}
Continux contributes four interrelated outcomes:

\begin{itemize}
\item A reproducible streaming lakehouse artifact deployed on a three-server K3s cluster using GitOps.
\item A Blue/Green materialized-view swap scenario measured while a public query loop runs continuously.
\item Observability dashboards with explicit interpretation guidance for cutover, resources, and data integrity.
\item An evidence-bounded analysis that separates documented platform capabilities from measured outcomes at this artifact.
\end{itemize}

Chapter 2 presents the theoretical background on lakehouse storage, stream processing, materialized views, GitOps, observability, and the Ursa paper. Chapter 3 describes the Design Science Research method, the Continux architecture, K3s topology, SQL objects, and Blue/Green design. Chapter 4 reports experimental results, including a source-data quality analysis, replay ingestion, cutover, and resource observations. Chapter 5 discusses results, limitations, and conclusions. The appendices preserve core SQL, configuration excerpts, evidence indexing, and a detailed terminology table.
